\chapter{Committing and executing}\label{ch:execute}

\section{Committing a schedule}\label{sec:commit}

When you are satisfied with an option, commit it. Committing says ``this is the plan''; work orders are
issued from the committed schedule, and its projected closing stock becomes the next period's opening stock.
Open the \ui{Commit} tab, or click \ui{Commit option} in the viewing bar.

\shot{commit}{The Commit tab: the option being committed, the mandatory reason and the decision log.}{commit}

\begin{steps}[Commit an option]
\begin{enumerate}
  \item Check \ui{You are committing}: the option's name and builder, on-time count, lateness, shortfall and
        the ranking policy. \ui{Change option} goes back to pick another one.
  \item Type the \ui{Reason} (required). It should say why this option and not another, for example ``accept
        3 extra shifts to keep enough F silver safety stock for long-dated orders''.
  \item Click \ui{Commit this option}. The status becomes \ui{Committed} and the next-step banner points to
        Execute.
\end{enumerate}
\end{steps}

The \textbf{decision log} below is stored on the server and is append-only: entries are never edited or
deleted. Each commit keeps the run, the data snapshot, the policy, the figures at that moment (late orders,
shortfall, machine shifts), the reason and who committed. You may commit another option until the period is
closed (\ui{Commit again}); the new row supersedes the previous one, and both remain visible.

\begin{note}
A decision must be explainable: the button refuses an empty reason (``Enter a reason before committing ---
the decision must be explainable''). If the server cannot be reached, the reason stays in the box so you can
try again.
\end{note}

\section{Work orders}\label{sec:work-orders}

The \ui{Execute} tab turns the committed schedule into work orders: what each machine does and when.

\shot{work_orders}{Work orders of Monday 23 November for all machines.}{work_orders}

\begin{itemize}
  \item Filter by \ui{Day} (or \ui{Whole period}) and by \ui{Machine} (or \ui{All machines}); the number of runs
        is shown.
  \item Each row gives the day and shift, the machine and stage, the run time, the run code and item, the
        quantity, and the preparation (when the changeover starts and how long it takes, the mold used, and a
        \ui{mold service} badge when maintenance precedes the run).
  \item \ui{CSV} exports the selection for Excel (with day, shift, machine, prep start, start, end, run, item,
        quantity, prep minutes and mold).
  \item \ui{Print} opens a printable sheet for the shift supervisors with a \ui{Done} column to tick and
        \ui{Prepared by} / \ui{Supervisor} signature lines.
\end{itemize}

If the period's inputs were changed after the commit, a banner warns that the work orders still follow the
committed schedule: if the change is real, generate new options and commit again.

\section{Incidents: replanning the rest of the period}\label{sec:incident}

When something goes wrong --- a machine breaks down, a rush order arrives, material is late --- you replan only
what has not happened yet. The \ui{An incident? Replan the rest} panel sits below the work orders:

\shot{incident_close}{Left: replanning after an incident. Right: finishing the period.}{incident_close}

\begin{steps}[Replan after a breakdown]
\begin{enumerate}
  \item Set the \ui{Cut-off}: the moment from which the plan changes (inside the period).
  \item Optionally choose the \ui{Machine down}, when it is \ui{Expected back} and \ui{What happened}. Leave
        \ui{None --- just replan} to replan without a breakdown (for example to add a rush order).
  \item Read the hint, which describes the state at the cut-off, for example ``By 30/11 08:00: 227/297 runs
        started, 21 orders delivered; still running on HT\_1 (until 16:00), HT\_2 (until 16:00)''.
  \item Click \ui{Replan from the cut-off}.
\end{enumerate}
\end{steps}

PlanWise creates a \term{reschedule} version of the period and opens its Prepare step:

\begin{itemize}
  \item what ran before the cut-off stays as it was (``fixed before'' the cut-off);
  \item runs in progress at the cut-off count as done, and their machine stays busy until they end;
  \item stock, molds and delivered orders are taken as of the cut-off; delivered orders leave the reschedule;
  \item the breakdown is entered as a calendar exception of the reschedule;
  \item pull-ahead is switched off (turn it back on under Policy if needed), so nothing is made ahead twice.
\end{itemize}

You can still adjust the reschedule (add a rush order, another stop), then generate options and commit as
usual. Section~\ref{sec:family} shows how reschedules appear among the versions of a period.

\section{Finishing the period}\label{sec:close}

The \ui{Finish the period} panel (right in Figure~\ref{fig:incident_close}) shows the projected closing stock
of each product against its safety level and how many orders the committed schedule delivers.

\begin{description}[style=nextline, leftmargin=1.2em, font=\bfseries\color{brand}]
  \item[Close only] Closes the period.
  \item[Close \& start the next period] Closes the period and creates the next one in one step, starting the
        day after with the projected closing stock as its opening stock. This is the normal end of a cycle.
\end{description}

Closing makes the period --- with its variants and reschedules --- read-only. The schedule that was executed
(the latest committed reschedule if there is one, otherwise the committed schedule of the period or variant)
is treated as what happened: the orders it delivered, including those delivered before each cut-off, are
marked \ui{Delivered} in the order book. Undelivered orders stay open and are picked up by the next period.

\shot{closed_period}{A closed period remains available for reference: its work orders, schedule, analysis and history.}{closed_period}

\begin{warning}
Close each period when it ends. A period that ended but is still open is flagged \ui{Overdue to close}, and the
orders it delivered remain open in the book, so a new period would take them again.
\end{warning}
